% !TEX root = ../journal_v2_0421052099.tex
%
% Condensed from thesis_book/chapters/introduction.tex.
\section{Introduction}\label{sec:intro}

\IEEEPARstart{L}{arge} language models answer factual questions fluently. But
often enough they give wrong answers, and those errors are not random. They are
systematic, and the model states them with the same confidence it gives a
correct answer~\cite{huang2025survey}. So a fabricated answer looks exactly like
a correct one. Things get worse when the answer requires putting several facts
together. A \emph{multi-hop} question such as ``Who directed the film that won
Best Picture in 1998?'' needs a chain of lookups in a fixed order. If the first
hop lands on the wrong entity, the rest of the chain carries on from a false
premise. And a check on the final answer alone never sees the intermediate
results.

Hallucinations arise at three points, and each calls for a different kind of
intervention~\cite{huang2025survey}. Defects in the training data can be fixed
only by curating the corpus and retraining. The model itself stores what it
knows diffusely across its weights, with no reliable signal that separates a
memorised fact from a plausible continuation. The third point is the prompt, and
it is the one this work takes on. When the context supplied with a question
already carries the facts needed to answer it, the model does not have to fall
back on parametric recall. Retrieval augmentation rests on that premise, and so
does this work. Our contribution is in how that context is assembled, and in
what gets checked before the answer goes out.

Retrieval-augmented generation (RAG)~\cite{lewis2020retrieval} attaches the
model to an external corpus and lets it answer from retrieved passages. For a
multi-hop question, the ranking signal measures the wrong thing. The passage
holding the intermediate entity may bear little resemblance to the question. And
flat retrieval cannot represent the relationships that a compositional question
has to traverse~\cite{li2023survey}. A knowledge graph makes those relationships
explicit, so composing two facts becomes a traversal.
GraphRAG~\cite{edge2024local} exploits this, but it still retrieves
\emph{statically}. It fetches a neighbourhood once, before reasoning starts. A
fixed radius is then either too small to contain the answer or so large that the
relevant edges are buried. The natural response is to interleave retrieval with
reasoning. Think-on-Graph~\cite{sun2024think} set that pattern with model-guided
beam search over the graph. Later systems added decomposition, backtracking, and
generated relation paths~\cite{chen2024plan,luo2024reasoning}.

Agentic navigation is the right shape for the problem, and this paper adopts it.
Three things it leaves unsolved motivate the work.

\textbf{Nothing checks what the final answer asserts.} Agentic systems ground
their answers in the subgraph they traverse, and the entities they assert do
exist in the graph. The problem is not fabrication but \emph{precision}. An
entity can sit in the retrieved subgraph, have been retrieved correctly, and
still be the wrong answer. Every system we survey emits whatever its final
generation call produces from the traversed context. Post-hoc self-checking
reduces factual error in free-form generation~\cite{dhuliawala2024chain}. What
none provides is a claim-level check against the retrieved triples, carried out
\emph{before} the answer is emitted, inside a graph-navigating agent.

\textbf{The contribution of individual mechanisms is unquantified.} Agentic
systems are assemblies of decomposition, scoring, backtracking, and
self-correction. The literature tends to report them as wholes. So when one
beats a static baseline, the mechanism that produced the gain is rarely
established.

\textbf{Accuracy is reported without cost.} Every exploration step costs at
least one model call. An accuracy table that leaves cost out makes agentic
methods look better across the board. The honest picture is a trade-off.

We present Agentic Graph Reasoning (AGR), a knowledge-graph question-answering
framework built as an explicit state machine over a constrained, deterministic
graph tool API. Its distinguishing component is a \emph{Structural Verification
Layer}. Before anything is emitted, the layer splits the draft answer into
atomic claims. It checks each one against the (head, relation, tail) triples the
agent actually traversed. Claims it cannot ground are re-explored or withdrawn.
So where it cannot ground a claim, AGR hedges rather than asserts. And the
answers it does give come paired with the triples that support them.

Stated formally, the input is a natural-language question $q$ and a knowledge
graph $G = (E, R, T)$ with entities $E$, relation types $R$, and triples $T
\subseteq E \times R \times E$. The system returns an answer set $A \subseteq E$
together with a supporting triple set $S \subseteq T$. The set $S$ contains only
triples the system traversed in reaching $A$, and it is checked before $A$ is
emitted. Three properties set this apart from standard question answering over
knowledge graphs. The output is a pair, and its supporting triples are part of
the contract. That is what makes the reasoning auditable rather than merely
asserted. The support must also be traversed, which rules out justifying, after
the fact, an answer the system reached some other way. And the justification is
checked before the answer is emitted. So unsupported content can be dropped or
hedged instead of merely flagged. Three research questions follow.

\begin{itemize}
\item \textbf{RQ1:} Does agentic navigation improve multi-hop factual accuracy,
  and does the advantage grow with the number of hops?
\item \textbf{RQ2:} What does pre-generation verification contribute beyond
  graph navigation?
\item \textbf{RQ3:} Which components contribute what, at what token cost?
\end{itemize}

We answer them with a controlled evaluation. All five systems share one frozen
backbone, one knowledge environment, and one call budget. So any difference
between them comes from architecture, not model capacity. The paper makes four
contributions.

\begin{enumerate}
\item AGR and its Structural Verification Layer (\Cref{sec:framework}). The word
  \emph{structural} names what the check establishes and, in the same word,
  bounds it. The check asks whether an edge joins a claim's endpoints, not
  whether it is the edge the claim names.
\item A controlled comparison of five systems on $400$ certified questions from
  each of WebQSP~\cite{yih2016value} and
  ComplexWebQuestions (CWQ)~\cite{talmor2018web}
  (\Cref{sec:setup,sec:results}). It
  splits the margin over the agentic baseline into search quality and
  affordability under a fixed budget.
\item A component-level ablation with paired significance tests
  (\Cref{sec:ablations}). It finds that
  decomposition is stratum dependent. Removing the planner improves accuracy on
  the shallower benchmark and trends toward harming it on the deeper one.
\item A hand reading of all $256$ remaining failures (\Cref{sec:errors}). It
  names the \emph{echo attractor}, a recurring failure in which several systems
  sharing a graph and a backbone converge on the same grounded but wrong answer.
  It also counts the benchmark's own label defects, which adjudication separates
  from those convergences.
\end{enumerate}

The headline results are these. AGR leads every accuracy column on both
datasets, and all eight paired tests against the baselines reject. Where the
shared call cap lets the agentic baseline finish, that baseline is slightly
ahead. So the margin is affordability, not better resolution. Structural
hallucination is eliminated by graph navigation itself, since the agentic
baseline also reaches zero. And removing the verification layer changes the
correctness of two questions out of $398$. Its contribution is auditability
rather than accuracy.

One limit applies to every absolute number. The knowledge environment is a union
of per-question subgraphs, each pre-extracted to contain its own question's
answer. That makes it friendlier than full
Freebase~\cite{bollacker2008freebase}, and the accuracies are not comparable
with evaluations over full Freebase. The comparisons between systems are sound,
because every system navigates the same environment (\Cref{sec:threats}).

The rest of the paper is organised as follows. \Cref{sec:related} reviews
related work, and \Cref{sec:environment} describes the knowledge environment.
\Cref{sec:framework} specifies AGR. \Cref{sec:setup} gives the experimental
setup and \Cref{sec:results} the results. \Cref{sec:errors} reads the failures.
\Cref{sec:discussion} discusses the findings and their limits, and
\Cref{sec:conclusion} concludes.

\endinput
